\section{Análise no Tempo Discreto}

\begin{frame}{Escolha do Período de Amostragem}
  \begin{columns}
    \begin{column}{0.62\textwidth}
      \begin{figure}
        \centering
        \includegraphics[width=\textwidth]{figuras/xA1_FR_fft.png}
        \caption{Espectro da resposta de $x_{A1}$ a um degrau em $F_R$.}
      \end{figure}
    \end{column}
    \begin{column}{0.34\textwidth}
      \begin{block}{Critério de Shannon}
        \[
          \omega_{\mathrm{amost}} \geq 2\omega_{\max}
        \]
      \end{block}
      \begin{block}{Período de amostragem}
        \[
          T_s = \frac{2\pi}{\omega_{\mathrm{amost}}}
        \]
      \end{block}
      \begin{block}{Valor adotado}
        \[
          \omega_{\mathrm{amost}} = 300 \ \mathrm{rad/h}
        \]
        \[
          T_s \approx 0{,}0209 \ \mathrm{h}
        \]
        \[
          T_s \approx 75{,}4 \ \mathrm{s}
        \]
      \end{block}
    \end{column}
  \end{columns}
\end{frame}

\begin{frame}{Resposta Contínua e Amostrada}

  \begin{figure}
    \centering
    \includegraphics[width=0.70\textwidth]{figuras/amostragem.png}
    \caption{Resposta contínua e amostrada de $x_{A1}$.}
  \end{figure}

  \vspace{-0.15cm}

  \begin{itemize}
    \item A resposta amostrada acompanha a resposta contínua.
    \item O período escolhido captura a variação inicial e a acomodação do sistema.
  \end{itemize}

\end{frame}

\begin{frame}{Estabilidade no Plano $z$}
  \begin{columns}
    \begin{column}{0.56\textwidth}
      \begin{figure}
        \centering
        \includegraphics[width=0.78\textwidth]{figuras/pzplot.png}
        \caption{Polos discretos no plano $z$.}
      \end{figure}
    \end{column}
    \begin{column}{0.40\textwidth}
      \begin{block}{Obtendo os polos discretos}
        $$z_i = e^{s_i T_s}$$
      \end{block}
      \begin{block}{Critério de estabilidade}
        $$|z_i|<1$$
      \end{block}
      \vspace{0.15cm}
      \begin{itemize}
        \item Todos os polos estão dentro do círculo unitário.
        \item A representação discreta permanece estável.
        \item O polo mais próximo do círculo unitário domina a resposta.
      \end{itemize}
    \end{column}
  \end{columns}
\end{frame}
